\chapter{El stack de \emph{crates}}
\label{cap:stack}

\begin{intuicion}
GraphHunter se divide en \emph{layers} concéntricos: \texttt{core}
(algoritmos puros), \texttt{platform} (contracts compartidos), y
\texttt{apps} (binarios end-user). La regla es que cada \emph{layer}
sólo importa de capas iguales o más internas: \texttt{core} no toca
\texttt{platform}, \texttt{platform} no toca \texttt{apps}. Esa
restricción aparenta arbitraria pero compra dos cosas concretas: el
\texttt{core} se puede testar en aislamiento (sin levantar Tauri,
sin levantar HTTP, sin Mock de Sentinel) y se puede reusar desde
varios \emph{transports} sin acoplamiento. Este capítulo dibuja el
diagrama, justifica el orden, y da una heurística práctica para
``¿en qué crate va esto que estoy escribiendo?''
\end{intuicion}

\section{Diagrama de capas}

\begin{figure}[H]
\centering
\begin{tikzpicture}[
    node distance=0.6cm,
    every node/.style={font=\footnotesize, align=center},
    layer/.style={draw=accent, thick, rounded corners=2mm,
                  minimum height=1cm, fill=accentlight},
    crate/.style={draw=accent!70, thick, rounded corners=1mm,
                  minimum height=0.7cm, font=\scriptsize\ttfamily,
                  fill=white}
]
\node[layer, minimum width=14cm] (apps) {\textbf{apps/}\,---\,binarios end-user};
\node[crate, below right=0.3cm and -7.5cm of apps.south west,
      minimum width=2.5cm] (cli)   {cli};
\node[crate, right=0.4cm of cli, minimum width=2.5cm] (tauri) {tauri};
\node[crate, right=0.4cm of tauri, minimum width=2.5cm] (gateway) {go-gateway};
\node[crate, right=0.4cm of gateway, minimum width=2.5cm] (mcp)   {mcp};

\node[layer, minimum width=14cm, below=2.0cm of apps] (plat)
    {\textbf{platform/}\,---\,API canónica + parsers + sources};
\node[crate, below right=0.3cm and -7.5cm of plat.south west,
      minimum width=2.0cm] (api)      {api};
\node[crate, right=0.3cm of api, minimum width=1.8cm] (dsl)      {dsl};
\node[crate, right=0.3cm of dsl, minimum width=2.0cm] (sources)  {sources};
\node[crate, right=0.3cm of sources, minimum width=1.8cm] (parsers) {parsers};
\node[crate, right=0.3cm of parsers, minimum width=2.0cm] (canonical) {canonical};
\node[crate, right=0.3cm of canonical, minimum width=1.6cm] (vrl)  {vrl};

\node[layer, minimum width=14cm, below=2.0cm of plat] (core)
    {\textbf{core/}\,---\,grafo temporal + matcher + scoring};
\node[crate, below right=0.3cm and -6.5cm of core.south west,
      minimum width=3.4cm] (engine)  {graph-engine};
\node[crate, right=0.4cm of engine, minimum width=3.0cm] (ffi)     {matcher-ffi};
\node[crate, right=0.4cm of ffi, minimum width=2.4cm] (gnn)        {gnn};

\draw[->, >=Latex, accent, very thick] (apps.south) -- (plat.north);
\draw[->, >=Latex, accent, very thick] (plat.south) -- (core.north);
\end{tikzpicture}
\caption{Layering. Las flechas son la dirección permitida de
\emph{imports}. \texttt{apps} importa de \texttt{platform} y
\texttt{core}; \texttt{platform} sólo de \texttt{core}; \texttt{core}
no ve hacia arriba. Excepciones puntuales se documentan en
\filepath{docs/arquitectura/layering-exceptions.txt} y obligan ADR.}
\label{fig:stack-layers}
\end{figure}

\section{Por qué este orden}

\begin{itemize}
    \item \textbf{\texttt{core}} es la pieza algorítmica pura. No
    conoce a Tauri, ni a HTTP, ni a Sentinel. Eso le permite testarse
    en aislamiento (\emph{property tests} con \texttt{proptest},
    \emph{benches} con \texttt{criterion}) y reusarse desde múltiples
    transports sin acoplamiento. Si mañana queremos exportar el
    matcher como \emph{Python binding}, lo único que cambiamos es
    una capa nueva en \texttt{platform}; el \texttt{core} no se
    toca.
    \item \textbf{\texttt{platform}} agrupa los \emph{contracts}
    compartidos entre transports: la \emph{canonical API}
    (\texttt{GraphHunterApi}), el parser de la DSL, los \emph{adapters}
    de \emph{sources} (Sentinel, Defender XDR, Data Lake), los
    parsers concretos (Cognito, FortiAnalyzer, IIS), y la
    \emph{canonical pipeline} (VRL + provenance). El criterio de
    inclusión: cualquier cosa que más de un transport puede querer.
    \item \textbf{\texttt{apps}} son los binarios usables: el
    \emph{daemon} Tauri (con su frontend React+Vite), el CLI que el
    \emph{Go gateway} invoca como sub-process, y el plugin MCP. Todos
    consumen \texttt{Arc<GraphHunterApi>} de \texttt{platform/api}.
\end{itemize}

\section{Heurística: ¿en qué crate va esto?}

Tres preguntas que guían la decisión cuando uno está escribiendo
código nuevo:

\begin{enumerate}
    \item \emph{¿Lo necesita más de un transport?}
    Sí $\Rightarrow$ \texttt{platform}. No $\Rightarrow$ \texttt{apps}.
    \item \emph{¿Tiene contrato externo (HTTP, MCP, FFI, REST)?}
    Sí $\Rightarrow$ siempre \texttt{platform}, nunca dentro de un
    \texttt{apps/} específico. Un \emph{transport} cualquiera puede
    querer reusar el contrato.
    \item \emph{¿Es algorítmico puro y libre de I/O o transport?}
    Sí $\Rightarrow$ \texttt{core}.
\end{enumerate}

\begin{ejemplo}
Tres dilemas reales y dónde fueron a parar:
\begin{itemize}
    \item ``Función que parsea un evento Sysmon NDJSON.''
    Pregunta 1: ¿más de un transport lo usa? Sí (Tauri y CLI ambos
    ingieren Sysmon). $\Rightarrow$ \texttt{platform/parsers}.
    \item ``Función que devuelve los vecinos salientes de un vértice
    ordenados por timestamp.'' Pregunta 3: ¿algorítmica pura? Sí.
    $\Rightarrow$ \texttt{core/graph-engine}.
    \item ``Componente React que renderiza la pildora ``Perf:M+L1''
    en el toolbar.'' Específico de Tauri. $\Rightarrow$
    \texttt{apps/tauri/src} (TypeScript).
    \item ``\emph{Endpoint} HTTP \texttt{POST /run\_hunt}.''
    Pregunta 2: contrato externo. $\Rightarrow$
    \texttt{platform/api} (la implementación), expuesto desde
    \texttt{apps/cli} cuando corre en modo daemon.
\end{itemize}
\end{ejemplo}

\section{Cómo se enforce la regla}

La capa \texttt{core} \emph{podría} importar de \texttt{platform} sin
que el compilador se queje (Cargo no impone direccionalidad). El
enforcement vive en CI:

\begin{enelcodigo}
\filepath{scripts/check-layering.sh} corre en cada CI run. Para cada
crate, lista sus dependencias declaradas en \texttt{Cargo.toml} y
verifica que no haya flechas ``hacia arriba''. Excepciones legítimas
(por ejemplo, un test que necesita un mock de \texttt{platform}) se
listan explícitamente en \filepath{docs/arquitectura/layering-exceptions.txt}
con justificación, y agregar una nueva entrada requiere ADR.
\end{enelcodigo}

Esa estrategia ---``layering por contrato + enforcement por
script + ADR para excepciones''--- es deliberadamente simple. Más
sofisticado (linter custom de \emph{cargo-deny} con reglas de
allowlist) sería más correcto pero menos legible. El script bash
gana.

\section{Crates en detalle}

\subsection{\texttt{core/}}

\begin{itemize}
    \item \texttt{graph-engine} — el grafo temporal in-memory,
    el matcher (DFS, LFTJ, Yannakakis), los índices auxiliares
    (NLF, k-hop), el scorer.
    \item \texttt{matcher-ffi} — la C ABI que expone el matcher para
    consumo desde C/C++/Python. Capítulo de la Especificación
    cubre la firma exacta.
    \item \texttt{gnn} — \emph{graph neural network} experimental
    para scoring; opcional (\emph{Cargo feature}).
\end{itemize}

\subsection{\texttt{platform/}}

\begin{itemize}
    \item \texttt{api} — la \texttt{GraphHunterApi}: el endpoint
    canónico que todos los transports consumen. \texttt{run\_hunt},
    \texttt{ingest\_batch}, \texttt{session\_save},
    \texttt{session\_load}, etc.
    \item \texttt{dsl} — el parser de la DSL (capítulo~\ref{cap:dsl-chain}).
    \item \texttt{sources} — adapters para Sentinel, Defender XDR,
    Data Lake, EVTX. Cada uno expone una \texttt{LogSource} trait
    que el ingestor consume.
    \item \texttt{parsers} — parsers de formato concretos: Cognito,
    FortiAnalyzer, IIS, Sysmon NDJSON, EVTX binario.
    \item \texttt{canonical} — la \emph{canonical pipeline} VRL +
    provenance: aplica un programa VRL a cada evento para
    normalizarlo a un esquema canónico antes de la ingesta.
    \item \texttt{vrl} — la integración con
    \emph{Vector Remap Language} para transformaciones declarativas.
\end{itemize}

\subsection{\texttt{apps/}}

\begin{itemize}
    \item \texttt{cli} — binario CLI que el Go gateway invoca como
    sub-process. Modo \emph{daemon} expone HTTP; modo \emph{one-shot}
    corre un comando y termina.
    \item \texttt{tauri} — el daemon con UI (React + Vite). El
    backend es Rust (Tauri commands), el frontend TypeScript. La
    aplicación que el analista usa día a día.
    \item \texttt{go-gateway} — un proxy Go que se sienta entre
    SIEM corporativo y el motor; traduce de la API del SIEM a la
    \texttt{GraphHunterApi}.
    \item \texttt{mcp} — plugin para el \emph{Model Context Protocol}
    que expone el motor como tool a un LLM.
\end{itemize}

\section{Por qué \emph{platform} y no \emph{infra} o \emph{shared}}

Una decisión de naming. ``Platform'' (en oposición a ``shared'',
``common'', ``utils'') refleja que esta capa NO es un cajón de
sastre sino una \emph{plataforma de transports}: cualquier nuevo
binario (futuro daemon Python, plugin Slack) construye encima.
\emph{Shared} y \emph{utils} se vuelven, por entropía, vertederos
de cualquier helper sin clasificar; \emph{platform} obliga a
preguntarse ``¿esto es realmente plataforma, o es código que
debería vivir junto al transport que lo usa?''

\begin{nota}
Si en algún punto uno se descubre escribiendo helpers genéricos sin
hogar (``utilidades de string'', ``wrappers para datetime''), la
respuesta correcta no es crear \texttt{platform/utils} sino dejarlos
en el módulo del consumer. Cuando dos consumers necesitan lo mismo,
\emph{ahí} se promueve a una crate compartida ---probablemente
fuera del repo, como dependencia de \texttt{Cargo.toml}.
\end{nota}

\section{Resumen}

\begin{itemize}
    \item Tres capas concéntricas: \texttt{core} (algorítmico puro),
    \texttt{platform} (contracts compartidos), \texttt{apps}
    (binarios). Imports sólo de afuera hacia adentro.
    \item Decidir en qué crate va algo: ¿más de un transport lo usa?
    ¿es contrato externo? ¿es algorítmico puro?
    \item Enforcement por \filepath{scripts/check-layering.sh} en CI;
    excepciones documentadas en \texttt{layering-exceptions.txt} y
    obligan ADR.
    \item El nombre \emph{platform} (no \emph{shared}, no
    \emph{utils}) es deliberado: evita que se vuelva basurero.
\end{itemize}

\begin{ejercicios}
\begin{enumerate}
    \item \starred{} Estás escribiendo una función que dado un
    \texttt{IP} string devuelve el ASN al que pertenece (vía
    look-up en una tabla MaxMind). ¿En qué crate va? Justificá con
    las tres preguntas guía.

    \item Una función que sólo usa el plugin MCP, que serializa una
    \texttt{Hypothesis} a un formato propio del MCP. ¿\texttt{platform}
    o \texttt{apps/mcp}? ¿Qué cambiaría si mañana también la
    quisiera usar el CLI?

    \item ¿Por qué la regla de layering NO se enforce en el
    compilador (Cargo) sino en CI (script)? Listá ventajas y
    desventajas de mover el check al compilador (vía
    \texttt{cargo-deny} u otro).

    \item \starred{} Imaginá que un \emph{hot path} del matcher
    necesita logging estructurado (cada operador emite eventos
    \texttt{tracing}). ¿Va en \texttt{core} (porque el operador
    está en \texttt{core}) o en \texttt{platform} (porque
    \texttt{tracing} es transport-aware)? Pensá los dos lados.

    \item Pensá en una \emph{exception} legítima: un caso donde sería
    razonable que \texttt{core} dependa de \texttt{platform}.
    ¿Qué pondrías en \texttt{layering-exceptions.txt} para
    documentarla?
\end{enumerate}
\end{ejercicios}
